\subsection{Thiết lập MLE: đi từ mật độ đến ước lượng}
\label{sec:mle_derivation}
Ước lượng hợp lý cực đại (maximum likelihood estimation, MLE) chọn các
tham số làm dữ liệu \emph{đã quan sát} có mật độ chung lớn nhất. Khi tìm
MLE, các $\mathbf x_j$ là số đã biết; $\vmu,\Sigma$ là các biến cần tối ưu.
Ta không tìm tham số có xác suất lớn nhất, vì tham số ở đây không được
đặt một phân phối xác suất.

Giả sử mẫu iid $\N_p(\vmu,\Sigma)$ và $\Sigma\succ0$. Để có nghiệm MLE
không suy biến, giả sử ma trận tổng tích ngoài quanh $\xbar$ xác định dương.
Trong mẫu chuẩn với hiệp phương sai tổng thể xác định dương, $n>p$ bảo
đảm điều này với xác suất 1; đối với bảng số cụ thể vẫn phải kiểm tra hạng.

\buoc{1}{Viết hàm hợp lý và log-hợp lý}
Tính độc lập cho phép nhân $n$ mật độ:
\begin{equation}
 L(\vmu,\Sigma)=(2\pi)^{-np/2}|\Sigma|^{-n/2}
 \exp\!\left[-\frac12\sum_{j=1}^n
 (\mathbf x_j-\vmu)'\Sigma^{-1}(\mathbf x_j-\vmu)\right].
 \label{eq:mlelikelihood}
\end{equation}
Logarit là hàm tăng, nên cực đại $L$ cũng là cực đại $\ell=\log L$:
\begin{equation}
 \ell(\vmu,\Sigma)=-\frac{np}2\log(2\pi)-\frac n2\log|\Sigma|
 -\frac12\sum_{j=1}^n(\mathbf x_j-\vmu)'\Sigma^{-1}(\mathbf x_j-\vmu).
 \label{eq:mleloglikelihood}
\end{equation}
Logarit chuyển tích thành tổng, làm phép tối ưu dễ xử lý hơn.
\nguon{Chương 4, §4.3, mật độ mẫu (4-11), tr.~168; hàm hợp lý (4-16), tr.~170}

\buoc{2}{Tách phần phụ thuộc trung bình và tìm $\widehat{\vmu}$}
Đặt
\[
 A=\sum_j(\mathbf x_j-\xbar)(\mathbf x_j-\xbar)'=(n-1)S.
\]
Vì $\sum_j(\mathbf x_j-\xbar)=0$, hai hạng chéo triệt tiêu khi khai triển:
\begin{equation}
 \sum_j(\mathbf x_j-\vmu)(\mathbf x_j-\vmu)'
 =A+n(\xbar-\vmu)(\xbar-\vmu)'.
 \label{eq:mlescatterdecomposition}
\end{equation}
Trace, ký hiệu $\tr(B)$, là tổng các phần tử đường chéo của ma trận vuông.
Với vectơ $\mathbf u$, ta có
$\mathbf u'B\mathbf u=\tr(B\mathbf u\mathbf u')$; do đó
\begin{align}
 \ell(\vmu,\Sigma)
 &=K-\frac n2\log|\Sigma|-\frac12\tr(\Sigma^{-1}A)
 \notag\\[-2pt]
 &\qquad-\frac n2(\xbar-\vmu)'\Sigma^{-1}(\xbar-\vmu),
 \label{eq:mlemeansplit}
\end{align}
trong đó $K=-np\log(2\pi)/2$ không phụ thuộc tham số.

Với $\Sigma$ cố định, chỉ hạng cuối phụ thuộc $\vmu$. Vì
$\Sigma^{-1}\succ0$, hạng toàn phương không âm, bằng 0 duy nhất khi
$\vmu=\xbar$. Nó mang dấu trừ trong log-hợp lý, nên
\begin{equation}
 \widehat{\vmu}=\xbar.
 \label{eq:mlemeanhat}
\end{equation}
Kết quả đúng với \emph{mọi} $\Sigma\succ0$ cố định, nên có thể thay
$\vmu=\xbar$ trước rồi tối ưu theo $\Sigma$.
\nguon{Chương 4, §4.3, phân rã (4-14), tr.~169; Result 4.11, tr.~171--172}

\buoc{3}{Đổi biến để tìm ma trận hiệp phương sai}
Đặt $\Theta=\Sigma^{-1}$, gọi là \emph{ma trận độ chính xác}. Đây là
phép đổi biến một-một trên các ma trận xác định dương. Vì
$\log|\Sigma|=-\log|\Theta|$, log-hợp lý sau khi chọn trung bình là
\begin{equation}
 \ell(\xbar,\Theta)=K+\frac n2\log|\Theta|-\frac12\tr(\Theta A).
 \label{eq:mleprecision}
\end{equation}
Ý nghĩa của hai hạng: $\log|\Theta|$ thưởng cho độ phân tán nhỏ, còn
$\tr(\Theta A)$ phạt nếu độ phân tán quá nhỏ so với các độ lệch thực tế.
Cực đại cân bằng hai hạng đó.

Để đạo hàm, dùng hai quy tắc ma trận
\[
 d\log|\Theta|=\tr(\Theta^{-1}d\Theta),\qquad
 d\tr(\Theta A)=\tr(A\,d\Theta).
\]
Ký hiệu $d\Theta$ là một thay đổi nhỏ của ma trận, không phải vectơ
tịnh tiến $\mathbf d$ dùng ở phần biến đổi dữ liệu. Hai quy tắc trên là
bản mở rộng của $(\log\theta)'=1/\theta$ và $(a\theta)'=a$.
Từ \eqref{eq:mleprecision},
\[
 d\ell=\frac12\tr\bigl[(n\Theta^{-1}-A)d\Theta\bigr].
\]
Tại điểm dừng, biểu thức phải bằng 0 theo mọi hướng thay đổi đối xứng,
nên $n\Theta^{-1}-A=0$. Suy ra
\begin{equation}
 \widehat\Sigma=\widehat\Theta^{-1}=\frac An
 =\frac1n\sum_j(\mathbf x_j-\xbar)(\mathbf x_j-\xbar)'
 =\frac{n-1}{n}S.
 \label{eq:mlecov}
\end{equation}
Đó là nguồn gốc của mẫu số $n$; ta không chọn mẫu số theo ý thích.

\buoc{4}{Kiểm tra điểm dừng thực sự là cực đại}
Đạo hàm bằng 0 chưa đủ để kết luận. Với hướng đối xứng $H\ne0$, đạo
hàm bậc hai theo đường $\Theta+tH$ là
\[
 \frac{d^2}{dt^2}\ell(\xbar,\Theta+tH)\bigg|_{t=0}
 =-\frac n2\tr(\Theta^{-1}H\Theta^{-1}H)<0.
\]
Thật vậy, trace trong vế phải bằng tổng bình phương các trị riêng của
ma trận đối xứng $\Theta^{-1/2}H\Theta^{-1/2}$, nên dương khi $H\ne0$.
Log-hợp lý là hàm lõm nghiêm ngặt theo $\Theta$ trên miền xác định dương,
vì vậy điểm dừng tìm được là cực đại duy nhất. Kết hợp Bước 2, đây cũng
là cực đại đồng thời theo cả trung bình và hiệp phương sai.
\nguon{Kết quả MLE: Chương 4, §4.3, Result 4.11, tr.~171--172.
 Sách chứng minh bằng Result 4.10; ở đây dùng dẫn xuất tương đương theo ma trận độ chính xác}

\hieu{Trung bình mẫu đặt tâm tốt nhất của đám mây. Hiệp phương sai MLE
khớp độ phân tán quanh tâm đó bằng trung bình của các tích ngoài.
Tối ưu theo ma trận độ chính xác giúp nhìn rõ nghiệm và kiểm tra cực đại.}

\subsection{MLE khác ước lượng không chệch như thế nào?}
Thay các giá trị quan sát bằng biến ngẫu nhiên, các công thức MLE trở
thành các \emph{ước lượng} $\widehat{\vmu}=\Xbar$ và
$\widehat\Sigma=(n-1)S/n$. Khi đã thay số, đó là các \emph{giá trị ước lượng}.
Ta có
\[
 \E(S)=\Sigma,\qquad
 \E(\widehat\Sigma)=\frac{n-1}{n}\Sigma.
\]
$S$ không chệch; MLE hiệp phương sai bị chệch xuống ở mẫu hữu hạn.
MLE tối ưu hàm hợp lý, không tối ưu tiêu chí ``không chệch''. Hai mục
tiêu khác nhau, nên cho mẫu số khác nhau. Khi $n$ lớn, tỉ lệ $(n-1)/n$
tiến về 1 và khác biệt giảm.

\begin{example}[Tính MLE từ mẫu bốn quan sát]
Với $\xbar=(6,10)'$, $S$ trong \eqref{eq:sample4}, $n=4$,
\[
 A=3S=\begin{pmatrix}24&-10\\-10&6\end{pmatrix},\quad
 \widehat{\vmu}=\begin{pmatrix}6\\10\end{pmatrix},\quad
 \widehat\Sigma=\frac A4=\begin{pmatrix}6&-2.5\\-2.5&1.5\end{pmatrix}.
\]
Ta có $|\widehat\Sigma|=2.75>0$ và $\widehat\Theta=4A^{-1}$,
kiểm tra được $4\widehat\Theta^{-1}=A$.
Để kiểm định Hotelling, ta dùng $S$, không thế $\widehat\Sigma$ vào
công thức $n\mathbf r'S^{-1}\mathbf r$ mà giữ nguyên ngưỡng $F$.
Để thiết lập tỷ số hợp lý, ta dùng đúng MLE $\widehat\Sigma$.
\end{example}

\luuy{Nếu $A$ suy biến, $A/n$ không thuộc miền $\Sigma\succ0$.
Không có MLE hữu hạn trong mô hình chuẩn không suy biến này: có thể
co phương sai về 0 theo một hướng không có độ lệch quan sát, khiến hàm
hợp lý tăng không giới hạn. Tình huống này xảy ra khi $n\le p$, hoặc
do các cột phụ thuộc tuyến tính. Phụ lục~\ref{app:mle_global} giải thích
bằng trị riêng, đồng thời kiểm tra cực đại mà không dùng đạo hàm ma trận.}

\kiemtra{Dùng cùng một mẫu, $S$ và MLE hiệp phương sai có thể khác nhau
không? Có; MLE bằng $(n-1)S/n$. MLE trung bình có cần biết trước
$\Sigma$ không? Không, nghiệm $\xbar$ đúng với mọi $\Sigma\succ0$.}
